% Bagian: computability
% Bab: recursive-functions
% Subbagian: introduction

\documentclass[../../../include/open-logic-section]{subfiles}

\begin{document}

\olfileid[id]{cmp}{rec}{int}
\olsection{Pendahuluan}

Untuk mengembangkan teori matematis tentang komputabilitas, pertama-tama kita
harus mengembangkan suatu \emph{model} komputabilitas. Sekarang kita memandang
komputabilitas sebagai jenis hal yang dilakukan oleh komputer, sedangkan
komputer bekerja dengan simbol. Namun, pada awal perkembangan teori
komputabilitas, contoh paradigmatik komputasi adalah komputasi
\emph{numerik}. Para matematikawan sejak dahulu tertarik pada fungsi-fungsi
teori bilangan, yakni fungsi $f\colon \Nat^n \to \Nat$ yang dapat dihitung.
Karena itu, tidak mengherankan bahwa pada awal teori komputabilitas,
fungsi-fungsi semacam itulah yang dipelajari. Contoh-contoh fungsi numerik
dapat dihitung yang paling dikenal, seperti penjumlahan, perkalian, dan
perpangkatan (bilangan asli), mempunyai suatu ciri yang menarik: semuanya
dapat didefinisikan secara \emph{rekursif}. Oleh sebab itu, cukup wajar untuk
mencoba memberikan definisi umum tentang \emph{fungsi dapat dihitung}
berdasarkan definisi rekursif. Di antara banyak cara yang mungkin untuk
mendefinisikan fungsi teori bilangan secara rekursif, satu pola definisi yang
sangat sederhana menjadi pokok bahasan utama di sini, yaitu apa yang disebut
\emph{rekursi primitif}.

Selain fungsi dapat dihitung, kita mungkin tertarik pada himpunan dan relasi
dapat dihitung. Suatu himpunan dapat dihitung jika kita dapat menghitung
jawaban atas pertanyaan apakah suatu bilangan yang diberikan merupakan
!!a{element} himpunan itu atau bukan, dan suatu relasi dapat dihitung jika dan
hanya jika kita dapat menghitung apakah suatu tupel
$\tuple{n_1, \dots, n_k}$ merupakan !!a{element} relasi itu atau bukan.
Dengan meninjau \emph{fungsi karakteristik} suatu himpunan atau relasi,
pembahasan himpunan dan relasi dapat dihitung dapat dicakup dalam pembahasan
fungsi dapat dihitung. Dengan demikian, kita juga dapat mendefinisikan relasi
rekursif primitif; misalnya, relasi ``$n$ membagi habis $m$'' merupakan relasi
rekursif primitif.

Namun, fungsi rekursif primitif---yakni fungsi yang dapat didefinisikan hanya
dengan rekursi primitif---bukanlah satu-satunya fungsi teori bilangan yang
dapat dihitung. Banyak perumuman rekursi primitif telah dipertimbangkan, tetapi
cara tambahan yang paling kuat dan diterima luas untuk menghitung fungsi ialah
pencarian tak berbatas. Cara ini menghasilkan definisi \emph{fungsi rekursif
  parsial} dan definisi yang berkaitan untuk \emph{fungsi rekursif umum}.
Fungsi rekursif umum dapat dihitung dan bersifat total, dan definisinya tepat
mencirikan fungsi rekursif parsial yang kebetulan bersifat total. Fungsi
rekursif dapat menyimulasikan setiap model komputasi lainnya (mesin Turing,
kalkulus lambda, dan sebagainya), sehingga merupakan salah satu dari banyak
model komputasi yang diterima.


\end{document}

